\section*{文档概览}

\BSRKeyPoint{4.14 / 方法校核}{从组织价值到可检测规约}{按《软件方法》上册第二版核对业务边界、执行者、消息方向与业务序列抽象级别；实物和纸单作为动作参数，不能作为有责任的业务生命线。重新区分前置状态与运行时验证，保留本地验证和真实平台边界。}

本文件是 BSR 滑雪租赁系统 4.14 版设计需求基线，以当前角色协作预览、008 签署付款顺序、009 预约安全重试、010 暂留到期可见性、011 本地备份演练、012 门店对象授权、013 本地启动安全闸门、014 本机候选身份和数据隔离、015 显式水平与重新备货释放、016 顾客确认换装与主数据一致性及 019 网页与小程序共享业务为依据更新。它面向课程设计、后续原型、数据库设计和编码验收，覆盖业务建模、执行者需求、系统用例、分析模型、模块化类图、数据模型、非功能要求和可追踪验收标准。附录 A～J 已纳入质量要求、验收场景、数据字典、迁移方案、追溯、参考资料及本轮规则，使单份 PDF 可以独立审阅。核心结论是：系统须按预约日期范围计算可用库存，以取货码连接线上订单与门店交付，逐件核验归还和押金；店员及经理只处理所属门店订单，经营负责人及管理员只看匿名汇总。异店接收凭完整代码核验，只显示未还实物和归还截止，不转让原门店费用结算责任。数据变更前须生成并验证恢复副本；课程预览只允许本机监听并防止公开目录读取私有数据。新预约须显式记录顾客本次水平；重新备货须释放旧占用并撤销旧检查。换装须由顾客本人确认准确差额后由门店执行，主数据写入与审计同成同败。微信小程序以独立会话经显式绑定访问同一顾客购物车、库存和订单；真实微信扫码与正式平台接口尚待验收。

\subsection*{本轮新增的协作闭环}
\BSRKeyPoint{04 / 人员与岗位}{每日排班对应实际到岗}{全部非顾客角色可记录本人上下班；经理看今天的岗位安排与打卡状态；已有出勤不能事后改班。}
\BSRKeyPoint{05 / 维修安全}{技师提交后仍须经理验收}{完整诊断、工时与四项安全检查后才进入待验收；通过才恢复库存，退回继续隔离，报废不可租。}
\BSRKeyPoint{06 / 角色联动}{同一数据库事实驱动各页面}{SSE 发变更信号，页面按自身权限重查；编辑时保留草稿，保存或关闭后同步。}
\BSRKeyPoint{07 / 到期透明}{旧单留痕，新单重新核价}{15 分钟暂留到期后显示失效、历史报价和重新选期入口；带回原装备但重新核价，库存释放，同键重试不延长暂留。}
\BSRKeyPoint{08 / 门店权限}{接收实物与查看订单分权}{店员和经理按所属门店访问逐单资料；异店凭完整代码仅办理逐件归还；经营和平台管理员只看匿名业务汇总。}
\BSRKeyPoint{09 / 本地预览}{启动与静态读取设安全闸门}{只允许本机监听；危险端口或数据库路径在打开数据库前拒绝；公开目录链接不能读取外部私有文件。}
\BSRKeyPoint{10 / 双端同源}{小程序与网页共用顾客订单}{微信身份经现有顾客账号证明后显式绑定；独立小程序会话只读写本人购物车与订单，同一库存和政策由服务端裁决。正式 AppID 已核对；真实登录、请求域名与扫码仍待验收。}
\subsection*{保留的三条核心设计红线}
\BSRKeyPoint{01 / 预约库存}{按日期承诺可用量}{按门店、装备规格和租期重叠范围计算库存；支付前重新校验，避免双重预订。}
\BSRKeyPoint{02 / 透明报价}{订单金额不可静默改变}{取货门店决定正式 NOK 价格；定位切换只影响参考展示币种，已确认订单保留价格快照。}
\BSRKeyPoint{03 / 押金闭环}{逐件验收后自动退款}{全部项目归还并通过验收且无争议，才自动发起押金原路退回；异常转人工审核。}
\subsection*{服务闭环}
\noindent\BSRFlowStep{01}{选择装备}{尺码·款式}\hfill\textcolor{Teal}{$\rightarrow$}\hfill
\BSRFlowStep{02}{预约支付}{日期·库存}\hfill\textcolor{Teal}{$\rightarrow$}\hfill
\BSRFlowStep{03}{到店取货}{取货码·验货}\hfill\textcolor{Teal}{$\rightarrow$}\hfill
\BSRFlowStep{04}{归还验收}{逐件检查}\hfill\textcolor{Teal}{$\rightarrow$}\hfill
\BSRFlowStep{05}{押金退回}{原路退款}
\subsection*{版本记录}

{\def\LTcaptype{none} % do not increment counter
\begin{longtable}[]{@{}
  >{\centering\arraybackslash}p{(\linewidth - 6\tabcolsep) * \real{0.0875}}
  >{\raggedright\arraybackslash}p{(\linewidth - 6\tabcolsep) * \real{0.1625}}
  >{\raggedright\arraybackslash}p{(\linewidth - 6\tabcolsep) * \real{0.6250}}
  >{\raggedright\arraybackslash}p{(\linewidth - 6\tabcolsep) * \real{0.1250}}@{}}
\toprule\noalign{}\rowcolor{TableHead}
\begin{minipage}[b]{\linewidth}\color{white}\bfseries\centering
\textbf{版本}
\end{minipage} & \begin{minipage}[b]{\linewidth}\color{white}\bfseries\centering
\textbf{日期}
\end{minipage} & \begin{minipage}[b]{\linewidth}\color{white}\bfseries\centering
\textbf{主要变化}
\end{minipage} & \begin{minipage}[b]{\linewidth}\color{white}\bfseries\centering
\textbf{状态}
\end{minipage} \\
\midrule\noalign{}
\endhead
\bottomrule\noalign{}
\endlastfoot
1.0 & 2026年10月 & 初版业务建模、现场租赁、维修与基础 ERD & 参考基线 \\
2.0 & 2026年10月20日 & 新增顾客在线预约、按日期库存、取货码、逐件验货、押金自动退款；补齐执行者需求、模块化类图、数据字典、验收和追踪 & 参考基线 \\
2.1 & 2026年10月21日 & 按课程总结补充需求获取与批准、用例规约字段、分析模型衔接、ER 映射、领域和表示层设计；改为六章目录 & 待 BSR 确认 \\
2.2 & 2026年10月22日 & 按 A---D 课程要求核实现状证据与蓝图边界，补齐执行者映射、逐用例序列图、可判定规约和类表一致性 & 待 BSR 确认 \\
2.3 & 2026年10月22日 & 补充挪威同业公开租赁价格调研、顾客端透明报价与多日价格规则；示例报价仅用于原型，正式金额待 BSR 批准 & 待 BSR 确认 \\
2.4 & 2026年10月23日 & 明确挪威市场价格以 NOK 计价，与国内人民币价格存在差异；顾客端显式标注币种 & 待 BSR 确认 \\
2.5 & 2026年10月23日 & 新增按定位切换参考币种与实时重算展示价；明确取货门店决定正式价、NOK 结算、汇率失效降级和订单金额锁定 & 待 BSR 确认 \\
3.0 & 2026年10月26日 & 按系统预览修订购物车与退出、全非顾客考勤、经理每日排班、维修档案及安全验收、跨角色同步；补齐用例、分模块类图、序列图、状态机、物理 ERD 和验收追踪 & 联动完善版 \\
4.0 & 2026年10月29日 & 统一 UC-01～15 与旧用例语义映射，重绘五对业务流程及模块化分析/数据图，标明现状证据和本地/目标差距 & 模型一致性版 \\
4.1 & 2026年10月29日 & 将 NFR、E-01～15、关键字段、迁移/恢复、源编号追溯和参考资料并回 PDF，恢复单份文档可独立审阅的完整性 & 待 BSR 确认 \\
4.2 & 2026年10月30日 & 顾客本人确认当前合同和现场纸质见证先于演示付款；合同发起人与经办人分离 & 本地已验证 \\
4.3 & 2026年11月2日 & 线上和现场开单按账号/渠道/提交标识去重，刷新找回原单；补齐规约、ERD、数据字典、验收与追溯 & 本地已验证 \\
4.4 & 2026年11月4日 & 暂留到期在列表、详情、流程与重试一致可见；旧单只读、库存释放并引导新租期，补齐状态图与验收 & 本地已验证 \\
4.5 & 2026年11月7日 & 在线生成一致性数据库副本、完整性与散列校验，并在隔离目录完成恢复演练；更新可恢复性需求与生产缺口 & 本地已验证 \\
4.6 & 2026年11月10日 & 店员与经理按门店授权；异店精确码最小资料接收；经营与管理员匿名汇总；补齐跨店写入和通知约束 & 本地已验证 \\
4.7 & 2026年11月11日 & 013 本地启动地址、端口和数据库路径闸门，公开静态链接逃逸拒绝；备份临时 sidecar 清理；71 项回归与隔离恢复复核 & 本地已验证 \\
4.8 & 2026年11月13日 & 按课件重绘业务序列图执行者、工人和实体图例；维修现状分开工人与装备；校正类图、分析序列图和 ERD 连线及文字间距，并逐页检查图示 & 图示已复核 \\
4.9 & 2026年11月17日 & 增补 014 本机候选的运行来源、离线首个管理员、HTTPS 会话、受控开通、隔离恢复及无演示数据边界；保留原六章和课堂图示 & 本机已验证 \\
4.10 & 2026年11月17日 & 014 独立审阅者 23 秒区分演示与候选边界；015 补充顾客本次水平显式选择、历史空值说明、重新备货释放旧单件和冲突回滚规约 & 本地验收版 \\
4.11 & 2026年11月17日 & 016 细化已交付订单的门店草案、顾客本人确认、纸质见证与演示收退；主数据业务写入和审计原子性、失败样本与历史兼容 & 规约补充 \\
4.12 & 2026年11月20日 & 019 增补微信小程序顾客渠道、独立身份绑定、共享购物车与报价预约、同单进度、模块类图/序列图/ERD；区分本机证据与真实平台验收 & 本机规约版 \\
4.13 & 2026年11月26日 & 核对正式 AppID 与小程序备案，修正原生页面、受控照片和验收状态；保留域名、真机与发布缺口 & 状态校核版 \\
4.14 & 2026年11月27日 & 按《软件方法》修正业务价值、组织边界、业务序列参与者及线型，区分实际交互与涉众，校核可检测前后置与签付备货顺序 & 方法校核版 \\
\end{longtable}
}

\subsection*{阅读与使用方式}

\begin{itemize}
\item
  课程设计报告可直接使用第 1 章至第 3 章说明项目依据、现状、角色和需求来源。
\item
  需求分析可使用第 3 章用例规约和第 4 章分析模型；用例路径中的主语只使用执行者名称或系统，避免混入页面按钮等交互细节。
\item
  系统设计可使用第 5 章和第 6 章的架构、模块化 ERD、数据字典、并发规则和权限矩阵。
\item
  开发与测试应以验收与追踪部分为入口，把需求编号映射到用例、类、表和验收场景；任何需求变更均应更新对应条目和版本记录。
\end{itemize}
